\section{A finite connection comparison}\label{sec:comparison}

The lattice argument will need to pass from a connection to a set of possible
relays to a connection to a target. We prove the finite inequality that makes
this passage possible, allowing different probabilities for the independent
bits so that the result remains available after some bits have been fixed.

A finite hypergraph has a finite node set and a finite labelled family of
nonempty subsets called hyperedges; distinct labels may represent the same
subset. Each hyperedge is independently open with its own probability; an
open hyperedge joins all its nodes. The cluster $C(x)$ always
contains $x$, and $C(S)=\bigcup_{x\in S}C(x)$. All connections in this section
are in finite undirected graphs or hypergraphs.

\begin{proposition}[Joint connection comparison]\label{cmp:joint}
In independent percolation on a finite hypergraph, with arbitrary hyperedge
probabilities in $[0,1]$, let $o$ be a node and let $A,T$ be node sets with
$A\ne\varnothing$. Then
\begin{equation}\label{cmp:joint-bound}
 \mathbb P(o\leftrightarrow A,\ o\leftrightarrow T)
 \ge \mathbb P(o\leftrightarrow A)
       \min_{a\in A}\mathbb P(a\leftrightarrow T).
\end{equation}
\end{proposition}

For bonds, subtraction immediately turns this inequality into the failure
bound needed later. Sites require an additional check because a hypergraph
node is always connected to itself, whereas a site self-connection requires
an open vertex. We give that check at the end of the section.

The proof of Proposition~\ref{cmp:joint} assigns each cluster meeting an
ordered list to its first node in that list. Conditional connection probabilities form the columns
of a triangular matrix. The positivity property below will bound the source's
joint connection probability from below by a nonnegative weighted sum of relay-to-target
probabilities. The weights will total the probability of reaching the relay
set, giving the proposition.

Deleting a node set $w$ means deleting its nodes and every hyperedge meeting
it; write $G\setminus w$ for the resulting hypergraph. On an exact-cluster
event $C(S)=w$ of positive probability, the hyperedges of $G\setminus w$
retain their independent original laws. Indeed, the event asks that every
node of $w$ connect to $S$ within $w$, and that all hyperedges crossing from
$w$ to its complement be closed. It tests no hyperedge disjoint from $w$.
The same observation applies after an earlier deletion.

\begin{lemma}[Cluster columns]\label{cmp:cone}
Let $G$ be a finite hypergraph whose hyperedges have independent open
probabilities strictly less than one.
Choose any ordered list of distinct nodes $x_1,\ldots,x_n$; the list need
not contain every node. Define
\[
 S_k=\{x_i:i<k\},\qquad
 E_k=\{C(x_k)\cap S_k=\varnothing\},\qquad
 d_k=\mathbb P(E_k)>0,
\]
and let $\mu_k$ be the conditional law of $C(x_k)$ given $E_k$, regarded
as a measure on
\[
 \mathcal D_k=\{U:x_k\in U,\ U\cap S_k=\varnothing\}.
\]
For a function $F$ on $\mathcal D_k$, let
\[
 v_k^F(i)=\mathbb E_{\mu_k}
       [\mathbf 1_{\{x_i\in U\}}F(U)],\qquad
 h_k=v_k^1,\qquad H=(h_1\ \cdots\ h_n).
\]
If $F$ is nonnegative and increasing under inclusion, then
\begin{equation}\label{cmp:cone-relation}
 v_k^F-\mu_k(F)h_k
   \in\operatorname{cone}\{h_j:k<j\le n\}.
\end{equation}
Here the cone consists of all nonnegative linear combinations, including
zero. In particular, $H^{-1}v_k^F$ is coordinatewise nonnegative.
\end{lemma}

\begin{proof}
Closing every hyperedge incident to $x_k$ proves $d_k>0$. The matrix $H$
is unit lower triangular: its $k$th column vanishes above row $k$ and has
entry one in row $k$. We prove \eqref{cmp:cone-relation} by descending
induction on $k$. For $k=n$, both $v_n^F$ and $\mu_n(F)h_n$ equal
$\mu_n(F)e_n$, where $e_n$ is the $n$th coordinate vector.

Fix $k<n$ and assume the assertion for every later column. Write
$s=x_k$, $S=S_k$, and $W=C(S)$. For a set $w\supseteq S$ avoiding $s$, put
\[
 c(w)=\mathbb E\bigl[F(C^{G\setminus w}(s))\bigr],
\]
where the expectation uses the original product law on the undeleted
hyperedges. Restriction of the same configuration shows that $c(w)$ is
nonincreasing in $w$.

\paragraph{A first-listed-node decomposition.}
Fix a row $i$. For $j>k$, let $L_{ij}$ be the event that the first listed
node of $C(x_i)$ is $x_j$.
Condition on an exact value $W=w$ avoiding $s$. Then the mean of $F(C(s))$
is $c(w)$. If $x_i\notin w$, its cluster contains the listed node $x_i$
and avoids all listed nodes preceding $s$, which belong to $w$.
Its first listed node is therefore $s$ precisely when $x_i\in C(s)$;
otherwise exactly one of the events $L_{ij}$ occurs.
On $L_{ij}$, condition further on $C(x_j)=U$. The set $U$ avoids $s$, and
the specification of $L_{ij}$ adds only membership and avoidance conditions
on $U$. Exact-cluster conditioning therefore leaves the mean of $F(C(s))$
equal to $c(w\cup U)$. Partitioning the expectation $c(w)$ and subtracting
$c(w)$ times the corresponding partition of total probability gives
\begin{align}
 &\mathbb E[\mathbf 1_{\{x_i\in C(s)\}}F(C(s))\mid W=w]
 \notag\\
 &\quad=c(w)\mathbb P(x_i\in C(s)\mid W=w)
   +\sum_{j>k}\mathbb E\bigl[
       \mathbf 1_{L_{ij}}\{c(w)-c(w\cup C(x_j))\}\mid W=w\bigr].
 \label{cmp:partition}
\end{align}
The expressions involving $c(w\cup C(x_j))$ are evaluated only on
$L_{ij}$. If $x_i\in w$, all $L_{ij}$ are empty under the conditioning,
and both sides vanish.

For $U$ disjoint from $S$, let
$W_U=C^{G\setminus U}(S)$. Define, for $U\in\mathcal D_k$,
\[
 (PF)(U)=\mathbb E[c(W_U)],
\]
and, for $U$ disjoint from $S\cup\{s\}$, define
\begin{equation}\label{cmp:residual}
 R_F(U)=\mathbb E\bigl[
   \mathbf 1_{\{s\notin W_U\}}
   \{c(W_U)-c(W_U\cup U)\}\bigr],
\end{equation}
with zero integrand when $s\in W_U$. The operator $P$ preserves
nonnegativity and monotonicity: enlarging $U$ shrinks $W_U$, whereas $c$
is nonincreasing.

We now average \eqref{cmp:partition} conditional on $E_k$. Given
$C(s)=U$ avoiding $S$, the conditional law of $W$ is that of $W_U$.
Thus the first term on the right averages to
\[
 \mathbb E_{\mu_k}[\mathbf 1_{\{x_i\in U\}}(PF)(U)]
 =v_k^{PF}(i).
\]
For a residual term, the exact event identity is
\[
 E_k\cap L_{ij}
 =E_k\cap E_j\cap\{x_i\in C(x_j)\}.
\]
Indeed, a cluster avoiding the earlier listed nodes also avoids their entire
clusters. Given $C(x_j)=U$ under $E_j$, the conditional law of $W$ is again
that of $W_U$, and the remaining restriction $E_k$ is $s\notin W_U$.
Consequently the unnormalized residual expectation is
\[
 \mathbb E\!\left[
  \mathbf 1_{E_k\cap L_{ij}}
   \{c(W)-c(W\cup C(x_j))\}\right]
 =d_j\,\mathbb E_{\mu_j}
       [\mathbf 1_{\{x_i\in U\}}R_F(U)],
\]
where the expression on the left is evaluated only when its indicator is one.
Dividing by $d_k$ and summing over $j>k$ proves the column identity
\begin{equation}\label{cmp:column-identity}
 v_k^F=v_k^{PF}+\sum_{j>k}\frac{d_j}{d_k}v_j^{R_F}.
\end{equation}
Here and below $R_F$ is restricted to $\mathcal D_j$ in the $j$th column.
To complete the induction, we show that $R_F$ is nonnegative and increasing
and that $P^tF$ converges uniformly to $\mu_k(F)$.

\paragraph{The residual is increasing.}
Sample all hyperedge bits once, and use their restrictions after every
deletion. For deterministic $U$ disjoint from $S\cup\{s\}$, the event
$W_U=w$ tests only bits of hyperedges disjoint from $U$ and meeting $w$.
Every full-graph hyperedge disjoint from $w$ therefore retains its independent
original law after this conditioning, including hyperedges that meet $U$.
When $s\notin w$, the two clusters in the following difference consequently
have the marginal laws defining $c(w)$ and $c(w\cup U)$.
Hence $R_F(U)=\mathbb E[D_U]$, where
\begin{equation}\label{cmp:residual-sample}
 D_U=\mathbf 1_{\{s\notin W_U\}}
 \left\{F(C^{G\setminus W_U}(s))
       -F(C^{G\setminus(W_U\cup U)}(s))\right\},
\end{equation}
again interpreted as zero on the excluded event. All clusters at which
$F$ is evaluated contain $s$ and avoid $S$.

The first cluster contains the second, so $D_U\ge0$. To compare two sets
$U\subseteq U'$, first note that $W_{U'}\subseteq W_U$. On
$\{s\notin W_U\}$, deleting the smaller set $W_{U'}$ makes the first
cluster larger. The second cluster has a different monotonicity:
\[
 C^{G\setminus(W_U\cup U)}(s)=C^{G\setminus U}(s).
\]
Indeed, in $G\setminus U$ the source cluster is disjoint from the union
$W_U$ of the $S$-clusters, so deleting $W_U$ removes none of its open
connections. The same equality holds for $U'$, because $s\notin W_{U'}$.
Thus the second cluster becomes smaller when $U$ is replaced by $U'$.
Both changes increase the difference of the two values of $F$, giving
$D_{U'}\ge D_U$. On $\{s\in W_U\}$, this inequality follows instead from
$D_U=0$ and $D_{U'}\ge0$. Taking expectations proves that $R_F$ is
nonnegative and increasing.

\paragraph{Conditional resampling and the limit.}
The operator $P$ is a Markov transition: from $U$, first sample $W_U$, then
independently sample $C^{G\setminus W_U}(s)$. Under the joint law of
$(C(s),W)$ conditional on $E_k$, these are the successive conditional laws
of $W$ given $C(s)$ and of $C(s)$ given $W$. Hence $\mu_k$ is stationary
for $P$. Alternating conditional cluster resampling under a separation
condition is the method used by van den Berg, H\"aggstr\"om, and
Kahn~\cite[Section~2.1, proof of Theorem~2.1]{BHK2006}; the cone conclusion here is proved by
the following iteration.

Every transition assigns probability at least
\[
 a=\prod_{e:\,s\in e}(1-p_e)>0
\]
to the singleton $\{s\}$, since closing every incident hyperedge suffices
to isolate $s$ in any deletion graph. Thus, for every real function $g$ on
the finite state space $\mathcal D_k$, subtracting this common mass gives
\[
 \max Pg-\min Pg\le(1-a)(\max g-\min g).
\]
Iteration makes the oscillation of $P^tF$ tend to zero. Its $\mu_k$-mean
remains $\mu_k(F)$ by stationarity, so
\[
 \sup_{U\in\mathcal D_k}|(P^tF)(U)-\mu_k(F)|
 \le(1-a)^t(\max F-\min F)\longrightarrow0.
\]
This proves the required uniform convergence, including at states of
$\mathcal D_k$ with zero $\mu_k$-mass.

By the induction hypothesis, each $v_j^{R_F}$ is a nonnegative combination
of $h_j,\ldots,h_n$. Equation~\eqref{cmp:column-identity} therefore puts
$v_k^F-v_k^{PF}$ in the cone on the right of
\eqref{cmp:cone-relation}. All iterates $P^tF$ remain nonnegative increasing,
so iteration puts $v_k^F-v_k^{P^tF}$ in that cone as well. The cone is closed,
because its generators belong to the basis given by the columns of $H$.
Taking the uniform limit proves \eqref{cmp:cone-relation}. Finally, adding
$\mu_k(F)h_k$ expresses $v_k^F$ as a nonnegative combination of columns of
$H$, proving the last assertion.
\end{proof}

\begin{proof}[Proof of Proposition~\ref{cmp:joint}]
First suppose all probabilities are below one. Use the notation of
Lemma~\ref{cmp:cone}. For
$F(U)=\mathbf 1_{\{x_j\in U\text{ for some }j>k\}}$, isolation of $x_k$
gives $\mu_k(F)<1$, and direct inspection of the rows gives
\[
 v_k^F-\mu_k(F)h_k=(1-\mu_k(F))(h_k-e_k),
\]
where $e_k$ is the $k$th coordinate vector. Lemma~\ref{cmp:cone} yields
nonnegative coefficients $\Gamma_{jk}$ such that
\[
 h_k=e_k+\sum_{j>k}\Gamma_{jk}h_j.
\]
With $\Gamma$ zero on and above the diagonal, these column identities say
$H=I+H\Gamma$, or $H^{-1}=I-\Gamma$. Also, for $d=(d_1,\ldots,d_n)^\mathsf T$,
\begin{equation}\label{cmp:first-node-mass}
 Hd=\mathbf 1,\qquad (I-\Gamma)\mathbf 1=d.
\end{equation}
Indeed, the cluster of $x_i$ has first listed node $x_k$ with probability
$d_kh_k(i)$, and these events partition the sample space.

If $o\in A$, \eqref{cmp:joint-bound} is immediate. Otherwise take the list
to consist exactly of the nodes of $A$, in any order, followed by $x_n=o$.
Set
\[
 F(U)=\mathbf 1_{\{U\cap A\ne\varnothing,\ U\cap T\ne\varnothing\}},
 \qquad y_i=\mathbb E[F(C(x_i))].
\]
Partitioning by the first listed node gives
$y=\sum_k d_kv_k^F$, where $F$ is restricted to $\mathcal D_k$ in each
summand. Every restriction is nonnegative increasing, so
Lemma~\ref{cmp:cone} gives $(I-\Gamma)y\ge0$ coordinatewise. In the last row,
\[
 y_n\ge\sum_{i<n}\Gamma_{ni}y_i.
\]
For $i<n$, $y_i=\mathbb P(x_i\leftrightarrow T)$, while
$y_n=\mathbb P(o\leftrightarrow A,\ o\leftrightarrow T)$. By
\eqref{cmp:first-node-mass}, the nonnegative weights in this row sum to
\[
 \sum_{i<n}\Gamma_{ni}=1-d_n=\mathbb P(o\leftrightarrow A).
\]
This proves \eqref{cmp:joint-bound}. Finally, finite event probabilities are
polynomials in the hyperedge parameters, and the minimum in
\eqref{cmp:joint-bound} is finite. Continuity therefore allows probabilities
equal to one as well.
\end{proof}

\begin{corollary}[Failure comparison]\label{cmp:failure}
Consider independent bond or site percolation on a finite undirected graph,
with arbitrary individual probabilities in $[0,1]$. For a vertex $o$, vertex
sets $A,T$, and $b\in[0,1]$, if
\[
 \mathbb P(a\leftrightarrow T)\ge1-b\qquad(a\in A),
\]
then
\begin{equation}\label{cmp:failure-bound}
 \mathbb P(o\leftrightarrow A,\ o\not\leftrightarrow T)\le b.
\end{equation}
The assertion therefore applies after any deterministic collection of bits
has been fixed.
\end{corollary}

\begin{proof}
The empty-set case $A=\varnothing$ is immediate. For bonds,
Proposition~\ref{cmp:joint} and subtraction give the stronger upper bound
$b\mathbb P(o\leftrightarrow A)\le b$.

For sites, use the incidence construction described
in~\cite[``One inequality for both models'']{SitePublic2026}.
Create a node $x^*$ for every graph vertex $x$ and a connector node for every
graph edge. For each site $x$, form a hyperedge consisting
of $x^*$ and all connector nodes corresponding to edges incident to $x$;
its independent open bit is the bit of site $x$. For distinct vertices
$x,z$, connection of $x^*$ to $z^*$ is equivalent to an open site path
from $x$ to $z$. Hypergraph self-connection, however, always holds.
Writing $A^*,T^*$ for the starred images, we consequently have
\[
 \mathbb P_{\mathrm{hyp}}(a^*\leftrightarrow T^*)
 \ge\mathbb P_{\mathrm{site}}(a\leftrightarrow T)\ge1-b.
\]
On the site event in \eqref{cmp:failure-bound}, the origin $o$ is open.
Its hypergraph connections to starred vertices therefore agree with its
site connections, including a possible zero-edge path to itself. This gives
the event inclusion
\[
 \{o\leftrightarrow A,\ o\not\leftrightarrow T\}_{\mathrm{site}}
 \subseteq
 \{o^*\leftrightarrow A^*,\ o^*\not\leftrightarrow T^*\}_{\mathrm{hyp}}.
\]
Applying Proposition~\ref{cmp:joint} and subtraction to the hypergraph bounds
the latter event by $b$, proving the site assertion. Fixed bits are exactly
parameters equal to zero or one, already allowed in the statement.
\end{proof}
